\documentclass[12]{article}
\usepackage[utf8]{inputenc}
\usepackage{authblk}
\usepackage{geometry}
\usepackage{amsmath}
\usepackage{hyperref}
\usepackage[french]{babel}
\usepackage{url}

\geometry{letterpaper,margin=2.5cm}

\hypersetup
{
    colorlinks = true, linkcolor = blue, citecolor = blue, urlcolor = blue,
}

%\def\labelitemi{$\bullet$}

\title{ECL 707 (Maîtrise) - ECL 807 (Doctorat) \\ Modélisation de la biodiversité}
%\date {23 juin 202}
\author{Dominique Gravel}
\affil{Département de biologie, Université de Sherbrooke}

\begin{document}

	\maketitle

	%-----------------------------
	\section*{Objectif général}

	La biodiversité est d'une remarquable complexité à différentes échelles spatiales. Comprendre les processus spatiaux responsables de cette variabilité est l'une des disciplines les plus actives, diversifiées et difficiles de l'écologie. Les populations, communautés et écosystèmes sont connectés dans l'espace par des mouvements d'individus, des flux de matière et d'énergie. Ainsi, comprendre la distribution de la biodiversité et prédire sa réponse aux changements globaux ne peut se faire sans considérer le voisinage et échanges entre écosystèmes. L'écologie spatiale aborde des thèmes aussi variés que la coexistence, l'utilisation de l'espace, l'organisation des réseaux trophiques, la stabilité, la mesure de la biodiversité et l'optimisation des aires protégées. L'objectif général de l'école d'été est d'explorer comment les processus spatiaux influencent l'organisation des communautés, du mouvement des individus jusqu'au fonctionnement des écosystèmes. Nous survolerons les théories de base et présenterons différentes techniques mathématiques, numériques et statistiques utilisées en écologie spatiale. L'enseignement combine des présentations magistrales, des exercices dirigés et la réalisation d'un travail d'équipe réparti sur l'ensemble de la semaine de formation. Des intervenants spécialisés sur différents thèmes participeront à la formation en offrant des séminaires sur des sujets et techniques de pointe. 

	L'école d'été s'adresse aux étudiants et étudiantes de deuxième et troisième cycle dans le domaine de la science de la biodiversité ainsi que dans les disciplines connexes telles que la foresterie, océanographie, limnologie et géomatique. 
	\vskip 1em    

	%-----------------------------
	\section*{Objectifs spécifiques}

	Au terme de ce cours, l'étudiant.e sera en mesure de:
    \vskip 1em 

	\begin{itemize}
	\renewcommand{\labelitemi}{$\bullet$}

		\item Conceptualiser un problème d'écologie spatiale en misant sur les principales théories écologiques

		\item Programmer différents types de modèles de dynamique spatiale

		\item Interpréter les propriétés émergentes 

		\item Analyser des données écologiques spatiales

		\item Identifier les limites les enjeux d'échelle spatiale pour une étude écologique

	\end{itemize}

    %-----------------------------

  \section*{Approche pédagogique}

  L'approche pédagogique repose sur de courtes leçons magistrales sur des notions théoriques de base, complétées d'exercices dirigés qui introduisent différentes techniques de modélisation spatiale. Des problèmes d'actualité en écologie spatiale seront soumis en début de semaine. Le travail se réalisera en équipe et les participant.e.s seront invité.e.s à présenter leurs résultats au terme de l'école intensive. 
  
	%-----------------------------
	\section*{Contenu}

    \begin{itemize}

        \item[Jour 1] Émergence de propriétés collectives à partir du mouvement individuel
		\item[Jour 2] Persistance sur des paysages complexes
		\item[Jour 3] Beta-diversité des métacommunautés compétitives
		\item[Jour 4] Biogéographie insulaire et écologie des réseaux trophiques
		\item[Jour 5] Interactions indirectes dans les méta-écosystèmes

    \end{itemize}

	%-----------------------------

	\section*{Évaluation}

	Les résultats du projet de programmation seront présentés et discutés en groupe à la dernière journée de l'école d'été. 
	L'évaluation portera sur la participation et une note finale de Succès ou Échec sera attribuée. 

\end{document}
